\section{Task, Data and Evaluation}
\label{sec:task}

\paragraph{Task.}
SemEval-2026 Task~6 \citep{thomas-etal-2026-semeval} uses the question--answer pairs of \citet{thomas-etal-2024-never}, taken from U.S.\ presidential interviews.
Each item is a triple: one sub-question (median 15 tokens), split out of the journalist's full turn (median 62 tokens), and the president's entire answer (median 266 tokens).
The label states how \emph{that sub-question} was answered.
Subtask~2 (\Stwo{}) has nine evasion types; Subtask~1 (\Sone{}) has three clarity levels that are a fixed function of them: Clear Reply (Explicit; 30.5\% of train), Ambivalent (Implicit, Dodging, General, Deflection, Partial/half-answer; 59.2\%) and Clear Non-Reply (Declining to answer, Claims ignorance, Clarification; 10.3\%).
Like the dataset paper and the two top-ranked systems \citep{shi-etal-2026-teleai,peerachaidecho-rutherford-2026-chulanlp}, we predict the 9-way label and derive the 3-way label from the taxonomy.

\paragraph{Data.}
Train has 3,448 items with one label each and an annotator identifier (three annotators, one per item).
Dev has 308 items labelled by three annotators: 125 are unanimous, 150 split 2--1 and 33 three ways, and 59.4\% carry more than one distinct label.
Test has 237 items with two annotators each; its labels were never released and the evaluation server has closed, so every comparison in this paper is on dev.
Two properties shape the modelling.
First, 69\% of training rows share their answer with another sub-question, and 71.7\% of those shared answers carry different labels, so the model must attribute parts of a long answer to the right sub-question.
Second, answers are long: for 101 of the 308 dev items the sub-question and answer exceed 512 DeBERTa tokens, and for 20 the full input exceeds 1,024.
Test answers are much shorter (median 71 tokens against 266 on train and 362 on dev).

\paragraph{Metric.}
Both subtasks are scored by macro-F1 over all declared classes; a class that is never predicted contributes $F_1=0$.
\Sone{} uses plain single-label macro-F1.
\Stwo{} uses \emph{multi-reference macro-F1}: with $G_i$ the set of annotator labels of item~$i$, a prediction $\hat{y}_i$ is a true positive for class~$\hat{y}_i$ if $\hat{y}_i \in G_i$, a false positive if not, and a miss charges one false negative to every member of $G_i$.
If every prediction is acceptable ($\hat{y}_i \in G_i$ for all $N$ items), there are no false positives or false negatives, every class predicted at least once has $F_1=1$, and
\begin{equation}
  \text{macro-F1} = \frac{\lvert\{\hat{y}_1,\dots,\hat{y}_N\}\rvert}{9} .
  \label{eq:task-coverage}
\end{equation}
Which acceptable label is named contributes nothing: a uniformly random member of each dev reference set scores 1.000 on five of five draws.
The score therefore rewards two things, landing in the reference set (the \emph{in-set rate}) and \emph{coverage}, naming each class at least once.
Since every class is worth one ninth of the score, a rare class such as Clarification, present in only 4 dev reference sets, counts as much as Explicit (115).
The floors in Table~\ref{tab:task-floors} show the trade-off: always predicting Explicit lands in-set for 37.3\% of items but scores 0.060, while uniform guessing lands in-set less often yet scores 0.129.

\begin{table}[t]
\centering
\small
\begin{tabular}{lccc}
\toprule
 & \multicolumn{2}{c}{\Stwo{}} & \\
\cmidrule(lr){2-3}
Floor (dev) & 3 ann. & 2 ann. & \Sone{} \\
\midrule
Always majority (Explicit) & 0.060 & 0.055 & 0.136 \\
Uniform random (20 draws) & 0.129 & 0.116 & 0.299 \\
TF-IDF + logistic regression & 0.257 & 0.243 & 0.445 \\
\bottomrule
\end{tabular}
\caption{Trivial floors on dev. ``2 ann.'': mean over the three two-annotator versions of the dev reference sets, the test regime. TF-IDF uses word 1--2-grams; its settings were chosen by 5-fold cross-validation on train.}
\label{tab:task-floors}
\end{table}

\paragraph{Evaluation protocol.}
The following protocol holds throughout the paper.
\emph{Epoch selection} uses the \emph{train slice}, a fixed stratified 10\% of train (345 rows, seed 12345) held out from training; dev is never used to choose anything.
Every system claim and every final comparison rests on \emph{10 paired seeds}, the same seed numbers for both arms, reported as mean $\pm$ standard deviation, with the per-seed difference, its $t$ statistic and the number of seeds that improve.
Variants are first screened on three seeds against a bar fixed in advance (+0.015 with at least two of three seeds up).
A \emph{single model} is one seed with argmax and no tuning on dev.
A \emph{system} averages the seeds' probabilities and applies logit adjustment \citep{menon2021longtail}:
\begin{equation}
  \hat{y}(x) = \arg\max_{c}\; \frac{\bar{p}(c \mid x)}{\pi_c^{\,\tau}},
  \label{eq:la}
\end{equation}
where $\bar{p}$ is the mean of the seeds' probabilities, $\pi_c$ the frequency of class~$c$ in train, and $\tau$ the only fitted parameter.
Because $\tau$ is fitted on dev, it is chosen by 5-fold nested cross-validation (grid $-1$ to $2$ in steps of 0.05): each fold's $\tau$ maximises \Stwo{} macro-F1 on the other four, and the pooled held-out predictions are scored.
We report the first cross-validation split and, where given, the mean over ten splits; \Sone{} is derived from the system's final 9-way prediction.
Two systems are compared by a bootstrap over dev items with the seeds fixed (2,000 resamples; \citealp{efron1993bootstrap}), which understates training-seed noise.
Because test has two annotators, we also rescore dev against each of its three two-annotator reference sets and report the mean.
Finally, the hypotheses and predicted outcomes of every experiment were recorded in the experiment log before it ran, and results are reported against them.
